\chapter{Implémentation et Déploiement}

\section{L'Écosystème Backend : Spring Boot et FastAPI}
L'implémentation de la logique métier robuste a été réalisée en Java 17, orchestrée par le puissant cadriciel Spring Boot 3. Les microservices ont été architecturés selon le modèle en couches, garantissant un découplage optimal et une testabilité accrue.

Parallèlement, la sphère de l'intelligence artificielle et du traitement d'image a été développée en Python en utilisant le cadriciel FastAPI. Ce choix technique est dicté par le support natif de la programmation asynchrone, indispensable pour gérer efficacement les délais de réponse incompressibles.

\section{L'Interface Utilisateur sous Angular}
L'expérience client a été forgée à l'aide d'Angular, le cadriciel front-end de Google. L'interface utilise les principes du Material Design pour offrir une ergonomie fluide. La programmation réactive s'est avérée fondamentale pour gérer l'affichage asynchrone des flux de réponse de l'intelligence artificielle et pour orchestrer les sondes de suivi comportemental.

\section{Conteneurisation et Orchestration avec Docker}
La standardisation du déploiement a été traitée avec la plus grande rigueur par l'adoption de la technologie Docker. L'ensemble des composantes applicatives, incluant les microservices Java, le backend Python, le client Angular, les instances PostgreSQL, la base vectorielle ChromaDB et le courtier Kafka, ont été encapsulées dans des conteneurs isolés et orchestrés par Docker Compose.
